\documentclass[11pt,a4paper]{article}
\usepackage[margin=2.4cm]{geometry}
\usepackage{amsmath,amssymb,amsthm}
\usepackage{booktabs}
\usepackage{enumitem}
\usepackage{hyperref}
\hypersetup{colorlinks=true,linkcolor=black,urlcolor=blue,
            pdftitle={How far the slice argument reaches}}

\newcommand{\Dset}{D}
\newcommand{\E}{E(P_{\Dset})}
\newcommand{\Z}{Z}
\newcommand{\R}{\mathbb{R}}
\newcommand{\Zint}{\mathbb{Z}}

\theoremstyle{plain}
\newtheorem{theorem}{Theorem}
\newtheorem{proposition}[theorem]{Proposition}
\newtheorem{corollary}[theorem]{Corollary}
\newtheorem{lemma}[theorem]{Lemma}
\theoremstyle{definition}
\newtheorem{definition}[theorem]{Definition}
\newtheorem{remark}[theorem]{Remark}
\newtheorem{example}[theorem]{Example}

\title{\bfseries How far the slice argument reaches\\[4pt]
\large A characterisation of the criteria for which the efficient set\\
of a multi-objective integer program is cheaply enumerable}
\author{}
\date{}

\begin{document}
\maketitle

\begin{abstract}
\noindent
Let $\Dset \subseteq \Zint^n_+$ be finite and let $\Z_1,\dots,\Z_p$ be criteria
on $\Dset$.  If $v$ is a non-dominated criterion vector then \emph{every}
$x \in \Dset$ with $\Z(x) = v$ is efficient, so the efficient set is exactly
the union of the \emph{slices} of the non-dominated set and no efficiency test
is needed on any slice point.  This argument assumes nothing about the
criteria.  What does need structure is the enumeration of a slice, and this
note determines exactly how much.

We prove: (i) a criterion has one-hyperplane level sets with data depending
affinely on the level \emph{if and only if} it is linear fractional, on the set
where the relevant denominator does not vanish --- so within that shape the
fractional case is maximal; (ii) the enumeration in fact needs only that each
level set be a finite union of \emph{polyhedra}, a strictly larger class that
contains $|c^{\top}x+a|/(d^{\top}x+b)$ and maxima and minima of ratios; and
(iii) a separate and strictly tighter condition governs the enumeration of the
non-dominated set itself, namely that sublevel sets be half-spaces.  The two
reaches do not coincide, and the smaller one is what confines such a method to
linear fractional criteria.  We close with what the characterisation implies
about neighbouring problem classes, and with a precise statement of what is
\emph{not} proved here.
\end{abstract}

\section{Setting}

Throughout, $\Dset \subseteq \Zint^{n}_{+}$ is a non-empty finite set --- in
the application $\Dset = \{x \in \Zint^n_+ : Ax \leqslant b\}$ with the system
bounded --- and $\Z = (\Z_1,\dots,\Z_p)$ is a vector of real-valued criteria on
$\Dset$, all to be maximised.

\begin{definition}
$y \in \Dset$ \emph{dominates} $x \in \Dset$ when $\Z(y) \geqslant \Z(x)$
componentwise with at least one strict inequality.  A point is \emph{efficient}
when nothing dominates it; $\E$ denotes the set of efficient points.  A vector
$v \in \R^p$ is \emph{non-dominated} when $v = \Z(x)$ for some $x \in \Dset$
and no $\Z(y)$ dominates it in the same sense; $F$ denotes the set of
non-dominated vectors.
\end{definition}

\begin{definition}
For $v \in \R^p$ the \emph{slice} at $v$ is
$S(v) = \{x \in \Dset : \Z(x) = v\}$.
\end{definition}

\section{The slice argument, which assumes nothing}

\begin{theorem}\label{th:slice}
If $v$ is non-dominated then every $x \in S(v)$ is efficient.
\end{theorem}

\begin{proof}
Let $x \in S(v)$ and suppose some $y \in \Dset$ dominates $x$.  Then
$\Z(y) \geqslant \Z(x) = v$ with at least one strict inequality, so $\Z(y)$
dominates $v$, contradicting the assumption that $v$ is non-dominated.  Hence
no $y$ dominates $x$.
\end{proof}

\begin{corollary}\label{cor:union}
$\displaystyle \E = \bigcup_{v \in F} S(v)$, and the union is disjoint.
\end{corollary}

\begin{proof}
($\supseteq$) is Theorem~\ref{th:slice}.  For ($\subseteq$), let $x$ be
efficient and put $v = \Z(x)$.  If some $\Z(y)$ dominated $v$ then $y$ would
dominate $x$, so $v \in F$ and $x \in S(v)$.  Disjointness holds because
$S(v) \cap S(w) = \emptyset$ for $v \neq w$.
\end{proof}

\begin{remark}
Theorem~\ref{th:slice} uses no property of $\Z$ whatsoever --- not linearity,
not continuity, not finiteness of $\Dset$ beyond what makes the notions
well posed.  So the \emph{characterisation} of $\E$ as a union of slices is
available for arbitrary criteria.  Everything that follows is about the cost of
\emph{enumerating} $S(v)$, which is a different question and the one where
structure enters.
\end{remark}

\begin{remark}\label{rem:cost}
The practical content of Corollary~\ref{cor:union} is that \textbf{no
efficiency test is paid on any slice point}.  A direct construction of $\E$
tests each candidate against the others; here the test is spent only in
producing $F$, and the points are recovered by enumeration.
\end{remark}

\section{Where one hyperplane is exactly the fractional case}

Call a criterion \emph{linear fractional} when
$\Z(x) = (c^{\top}x + a)/(d^{\top}x + b)$ for fixed $c,d \in \R^n$ and
$a,b \in \R$.

\begin{proposition}\label{prop:lf-gives-hyperplane}
If $\Z$ is linear fractional and $d^{\top}x + b \neq 0$ on $\Dset$, then for
every level $v$,
\[
  S(v) = \{x \in \Dset : (c - v d)^{\top} x = v b - a\} ,
\]
an intersection of $\Dset$ with one hyperplane whose data is affine in $v$.
\end{proposition}

\begin{proof}
$\Z(x) = v$ iff $c^{\top}x + a = v(d^{\top}x + b)$, since multiplying by the
non-vanishing $d^{\top}x+b$ is reversible; rearranging gives the stated
equation, whose coefficient vector $c - vd$ and right-hand side $vb-a$ are
affine in $v$.
\end{proof}

The converse is the part that makes this a characterisation, and it holds for
\emph{any} affine dependence on the level rather than only for the
parametrisation just produced.

\begin{theorem}[Limit]\label{th:limit}
Suppose there are $a_0, a_1 \in \R^n$ and $b_0, b_1 \in \R$ with
\[
  \{x : \Z(x) = v\} = \{x : (a_0 + v a_1)^{\top} x = b_0 + v b_1\}
  \qquad \text{for every level } v .
\]
Then on $\{x : b_1 - a_1^{\top}x \neq 0\}$ the criterion is linear fractional:
\[
  \Z(x) \;=\; \frac{a_0^{\top}x - b_0}{\,b_1 - a_1^{\top}x\,} .
\]
Together with Proposition~\ref{prop:lf-gives-hyperplane}: \textbf{on the set
where the denominator does not vanish, a criterion has one-hyperplane level
sets with data affine in the level if and only if it is linear fractional.}
\end{theorem}

\begin{proof}
Fix such an $x$ and put $v = \Z(x)$.  Then $x$ lies in its own level set, so
$(a_0 + v a_1)^{\top}x = b_0 + v b_1$, that is
$a_0^{\top}x - b_0 = v\,(b_1 - a_1^{\top}x)$.  Dividing by the non-zero
$b_1 - a_1^{\top}x$ gives the claim, with $c = a_0$, $a = -b_0$, $d = -a_1$
and $b = b_1$.
\end{proof}

\begin{remark}
The non-vanishing hypothesis is not a technicality to be dropped.  It is the
positive-denominator assumption under which the fractional problem is normally
posed, restated; without it $\Z$ is not even continuous and the level-set
description can fail.  The theorem is accordingly stated \emph{on} that set.
\end{remark}

\begin{remark}
Theorem~\ref{th:limit} is a \emph{limitation}, and that is its value.  It says
there is nothing larger \emph{of that shape} to look for: a method whose slices
are single hyperplanes with affine data is a method for linear fractional
criteria, whatever else one might hope.
\end{remark}

\section{But one hyperplane was never required}

What an enumerator of $S(v)$ actually consumes is a system of linear
constraints.  Nothing forces that system to describe a single hyperplane, and
nothing forces it to be a single system.

\begin{proposition}[Reach, affine version]\label{prop:reach-affine}
Suppose that for each level $v$ there are finitely many affine subspaces
$A_1(v),\dots,A_{m}(v)$, computable from $v$, with
$\{x : \Z(x) = v\} = \bigcup_{i} A_i(v)$.  Then
$S(v) = \bigcup_i \bigl(\Dset \cap A_i(v)\bigr)$ is enumerable at $m$ times the
cost of the one-subspace case.
\end{proposition}

\begin{proof}
Immediate: enumerate $\Dset \cap A_i(v)$ for each $i$ and take the union,
discarding repetitions where the $A_i(v)$ meet.
\end{proof}

\begin{example}\label{ex:abs}
$\Z(x) = |c^{\top}x + a| / (d^{\top}x + b)$ with $d^{\top}x + b > 0$ is not
linear fractional, and
\[
  \{x : \Z(x) = v\} = \{x : c^{\top}x + a = v(d^{\top}x+b)\}
  \;\cup\; \{x : c^{\top}x + a = -v(d^{\top}x+b)\} ,
\]
two hyperplanes.  By Proposition~\ref{prop:reach-affine} its slices are
enumerable.
\end{example}

The same argument runs once more, and the second step is the one that reaches a
criterion type of independent interest.

\begin{proposition}[Reach, polyhedral version]\label{prop:reach-poly}
Suppose that for each level $v$ there are finitely many polyhedra
$P_1(v),\dots,P_m(v)$, computable from $v$, with
$\{x : \Z(x) = v\} = \bigcup_i P_i(v)$.  Then $S(v)$ is enumerable at $m$ times
the cost of one polyhedron.
\end{proposition}

\begin{proof}
As before; an enumerator that accepts linear equations accepts linear
inequalities at no structural cost, since both are linear constraints on the
same variables.
\end{proof}

\begin{example}[Minimax criteria]\label{ex:max}
Let $\Z(x) = \max_{i \leqslant m} \Z_i(x)$ with each $\Z_i$ linear fractional
of positive denominator.  Then
\[
  \{x : \Z(x) = v\} \;=\; \bigcup_{j=1}^{m}
  \Bigl\{ x : \Z_j(x) = v \ \text{ and } \ \Z_i(x) \leqslant v
  \ \ (i \neq j) \Bigr\} ,
\]
and each branch is one linear equation together with $m-1$ linear inequalities
--- a polyhedron, not an affine subspace.  Minimisation is the same with the
inequalities reversed.
\end{example}

\begin{remark}
Propositions~\ref{prop:reach-affine} and~\ref{prop:reach-poly} are one
observation applied twice: what a slice enumeration needs is not a \emph{shape}
but \emph{finitely many linear systems}.  What it cannot absorb is curvature.
\end{remark}

\section{The other half of such a method reaches less}

Corollary~\ref{cor:union} presupposes $F$.  Producing $F$ is a separate
problem, and a method that searches criterion space by splitting around a
visited point faces a different requirement.

Such a search, at a centre $\bar x$, divides the remaining region into
$\{x : \Z_k(x) > \Z_k(\bar x)\}$ and $\{x : \Z_k(x) \leqslant \Z_k(\bar x)\}$.
For each part to be a single linear program, both must be half-spaces.

\begin{proposition}[Necessary condition for the split]\label{prop:front}
If the split above is to express each part as one linear program for every
centre, then for every \emph{attained} level $t \in \Z(\Dset)$ the sublevel set
$\{x : \Z(x) \leqslant t\}$ must be a half-space.
\end{proposition}

\begin{proof}
Taking $t = \Z(\bar x)$ as $\bar x$ ranges over $\Dset$, the sublevel set is
one of the two parts, hence a half-space; the other part is its complement,
which is then a half-space as well.
\end{proof}

\begin{remark}
The restriction to attained levels is not cosmetic: the split only ever
produces levels of the form $\Z(\bar x)$, so that is all the argument delivers.
A statement about every real $t$ would be stronger than the proof, and on a
finite $\Dset$ the attained levels are all that the enumeration ever sees.
\end{remark}

\begin{remark}[What is not claimed]\label{rem:notexact}
Proposition~\ref{prop:front} is a \emph{necessary} condition, not a
characterisation.  A criterion could satisfy it and still not be usable,
because a method must also \emph{construct} those half-spaces from the problem
data, and nothing above provides that in general.  Sufficiency is available for
linear fractional criteria, where after scaling the data to integers the
quantity
\[
  e_k(x ; \bar x) \;=\; D_k(\bar x)\,D_k(x)\,
  \bigl(\Z_k(x) - \Z_k(\bar x)\bigr)
\]
is linear in $x$ and integer-valued, so that
$\Z_k(x) > \Z_k(\bar x) \iff e_k(x;\bar x) \geqslant 1$ exactly and both parts
are explicit.  An earlier draft of this material asserted an equivalence here;
that assertion is withdrawn.
\end{remark}

\begin{proposition}\label{prop:gap}
The condition of Proposition~\ref{prop:front} is strictly stronger than that of
Proposition~\ref{prop:reach-poly}.
\end{proposition}

\begin{proof}
For the containment, let $\Z$ satisfy Proposition~\ref{prop:front} and let
$t_1 < \dots < t_N$ enumerate $\Z(\Dset)$, which is finite.  Write
$H_i = \{x : \Z(x) \leqslant t_i\}$, a half-space by hypothesis.  Since the
levels below $t_i$ are exactly $t_1,\dots,t_{i-1}$,
\[
  \{x : \Z(x) = t_i\} \;=\; H_i \setminus H_{i-1}
  \qquad (H_0 = \emptyset) ,
\]
an intersection of a half-space with the complement of a half-space, hence a
polyhedron.  So Proposition~\ref{prop:reach-poly} applies with $m = 1$.

\emph{(The step one is tempted to write --- that the boundary of a half-space
sublevel set is the level set --- does not hold in general, and on a finite
$\Dset$ a topological boundary has no content.  The difference of consecutive
sublevel sets is what carries the argument.)}

Strictness: in Example~\ref{ex:abs} the
sublevel set $\{|c^{\top}x+a| \leqslant v(d^{\top}x+b)\}$ is the intersection
of two half-spaces and its complement is not convex; in Example~\ref{ex:max} it
is $\bigcap_i \{\Z_i \leqslant v\}$, a polyhedron whose complement is a union.
Neither is a half-space, while both have enumerable slices.
\end{proof}

\begin{remark}
Proposition~\ref{prop:gap} explains a feature of such methods that is otherwise
puzzling.  They are described as methods for linear fractional programs, yet
their completion step is indifferent to the criteria.  The reason is that the
two halves have different reaches, and \textbf{the smaller one names the
method}.  For a criterion between the two --- Example~\ref{ex:max} is one ---
the efficient set can be completed exactly once $F$ is known, but $F$ must be
obtained some other way.
\end{remark}

\section{A consequence for the decision function}

\begin{theorem}[Admissible decision functions]\label{th:admissible}
Let $\E$ be known.  Then for any function $f$ defined on $\Dset$, both
$\max\{f(x) : x \in \E\}$ and $\operatorname*{argmax}\{f(x) : x \in \E\}$ are
obtained by evaluating $f$ on $\E$.  No linearity, monotonicity, continuity,
differentiability or ratio form is required of $f$.  The same holds for the
non-dominated subset of $\E$ under finitely many such functions.
\end{theorem}

\begin{proof}
$\E$ is finite.
\end{proof}

\begin{remark}
The proof is the whole of the content, and stating the theorem would be idle
were it not for Corollary~\ref{cor:union} and Remark~\ref{rem:cost}: the point
is not that maximising over a finite set is easy, but that the slice argument
makes $\E$ affordable, and that this is what the affordability buys.  Methods
that never construct $\E$ must instead derive their cuts and bounds from the
form of $f$, and so are confined to the forms they can differentiate or
linearise.
\end{remark}

\section{What the characterisation says about neighbouring classes}

Theorem~\ref{th:limit} and Proposition~\ref{prop:reach-poly} classify a
criterion by inspection of its level sets.

\begin{example}[Sums of a linear and a fractional term]\label{ex:lplf}
Let $\Z(x) = c^{\top}x + (p^{\top}x + \alpha)/(q^{\top}x + \beta)$.  Clearing
the denominator at level $v$,
\[
  (c^{\top}x)(q^{\top}x + \beta) + p^{\top}x + \alpha
     - v\,(q^{\top}x + \beta) \;=\; 0 ,
\]
whose quadratic part is $(c^{\top}x)(q^{\top}x)$, with matrix
$(c q^{\top} + q c^{\top})/2$.  That matrix is non-zero whenever $c$ and $q$
both are, so the level set is a conic.  A conic is not a finite union of
polyhedra, so Proposition~\ref{prop:reach-poly} does not apply and the slice
machinery does not reach this class.
\end{example}

\noindent
The same computation disposes of criteria that are products, and of quadratic
criteria, for the same reason: curvature.  Example~\ref{ex:lplf} is worth
isolating because the class it describes is studied in its own right, so the
characterisation says in advance that a slice-based method and a method for
that class are addressing disjoint problems.

\section{What is not proved here}

\begin{enumerate}[itemsep=3pt]
\item \textbf{Proposition~\ref{prop:front} is not a characterisation.}  See
      Remark~\ref{rem:notexact}.  Sufficiency is shown only for linear
      fractional criteria.
\item \textbf{Nothing here bounds the cost.}  Propositions~\ref{prop:reach-affine}
      and~\ref{prop:reach-poly} say a slice is enumerable, not that it is
      small.  The size of $S(v)$ and the number of branches are properties of
      the instance.
\item \textbf{Theorem~\ref{th:limit} concerns level sets with data affine in
      the level.}  A criterion whose level sets are single hyperplanes with
      data depending on $v$ in some other way is not covered, and we do not
      claim such criteria cannot exist.
\item \textbf{The classification of a published problem class} --- as in
      Example~\ref{ex:lplf} --- rests on the class as stated.  The algebra is
      ours; the identification of which class a given paper treats is not a
      mathematical claim.
\end{enumerate}

\end{document}
